Geometric compactness metrics---PP, Reock, CH, S, LW---measure properties
of district shapes as geometric objects.
They answer the question: \emph{Is this shape compact?}
Population-Weighted Compactness (PWC) answers a different question:
\emph{Are the residents of this district well-served by their geographic arrangement?}

A district can be perfectly circular (PP $= 1$, Reock $= 1$, CH $= 1$)
but have all its population concentrated at the district's periphery,
forcing residents to live far from the district centroid and from each other.
Conversely, a geometrically irregular district can have low PWC if
its population is concentrated near a central urban core.

\subsection{The Fryer-Holden Formulation}

\citet{fryer2011measuring} introduced a moment-of-inertia compactness metric
equivalent to PWC and showed that it captures a dimension of representational
quality distinct from geometric metrics.
In their framework, PWC represents the expected squared distance between two
randomly drawn residents of the same district---a measure of how ``close together''
the district's population is.
\citet{fryer2011measuring} find that minimising PWC is equivalent to minimising
travel costs for district-level political activity (town halls, canvassing,
constituent service).

\subsection{Relationship to Community of Interest}

Many state redistricting guidelines include a ``community of interest'' criterion:
districts should contain ``geographically compact'' communities with shared
interests~\citep{ncsl2021criteria}.
PWC provides a data-driven operationalisation of this criterion for the spatial
component: communities with low PWC have their residents clustered near a
common geographic centre, facilitating the shared-interest representation
that the community-of-interest standard seeks to protect.

\subsection{Scope}

This paper provides the first systematic comparison of PWC across all four
bisect structure algorithms on NC, WI, and TX (2020 Census, seed $= 42^\dagger$).
We document the prime-factor algorithm's PWC advantage and the low correlation
between PWC and geometric metrics.
